\documentclass[a4paper]{article}
\usepackage{amsfonts}
\usepackage{amsmath}
\usepackage{amssymb}
\usepackage{graphicx}     

\newcommand{\dint}{\displaystyle\int}
\newcommand{\llim}{\lim\limits}
\newcommand{\Bgtan}{\operatorname{Bgtan}}

\begin{document}
  
\noindent \large Auteur: Oubayy Ahale         \\
\noindent \large Studierichting: Informatica  \\
\noindent \large Oefeningenreeks 5A nr. 5.6   \\

\medskip

\normalsize

Bereken de integraal:

$\dint_{0}^{\tfrac{\pi}{2}} \sin^2(x) dx$\\

$\dint_{0}^{\tfrac{\pi}{2}} \sin^2(x) dx = \dint_{0}^{\tfrac{\pi}{2}} \dfrac{1 - \cos(2x)}{2} dx = \dfrac{1}{2} \cdot \dint_{0}^{\tfrac{\pi}{2}} \left( 1 - \cos(2x) \right) dx$\\

$= \dfrac{1}{2} \cdot \left[ x - \dfrac{\sin(2x)}{2} \right]_{0}^{\tfrac{\pi}{2}}$\\

Voor $x = \dfrac{\pi}{2}$:\\

$x - \dfrac{\sin(2x)}{2} = \dfrac{\pi}{2} - \dfrac{\sin(\pi)}{2} = \dfrac{\pi}{2} - 0 = \dfrac{\pi}{2}$\\

Voor $x = 0$:\\

$x - \dfrac{\sin(2x)}{2} = 0 - \dfrac{\sin(0)}{2} = 0$\\

Dus:\\

$\dfrac{1}{2} \cdot \left( \dfrac{\pi}{2} - 0 \right) = \dfrac{\pi}{4}$\\

\textbf{Antwoord:} $\boxed{\dfrac{\pi}{4}}$\\

\end{document}

\begin{thebibliography}{999}
%\bibitem{a} Referentie
\end{thebibliography}
\end{document}